\begin{corollary}[5.12.10]
\label{IV.2.s5.12.10-fr}
Sea $A$ un anillo local noetheriano íntegro que verifica una de las hipótesis
siguientes:
\begin{enumerate}
\item[\textup{a)}] $A$ es cociente de un anillo regular.
\item[\textup{b)}] $A$ es universalmente catenario y universalmente japonés.
\end{enumerate}
Sea $(x_i)_{1\leq i\leq n}$ una sucesión de $n$ elementos de $A$; pongamos
$\mathfrak J=x_1A+\ldots+x_nA$, y supongamos que se verifican las condiciones
siguientes:
\begin{enumerate}
\item[\textup{(i)}] Solo existe un ideal primo asociado no inmerso $\mathfrak p$ de
$A/\mathfrak J$, y $\mathfrak p$ es de altura $n$.
\item[\textup{(ii)}] El ideal maximal $\mathfrak pA_{\mathfrak p}$ de
$A_{\mathfrak p}$ está engendrado por los $x_i/1$ \textup{(por tanto,
$A_{\mathfrak p}$ es un anillo local \emph{regular} de dimensión $n$)}.
\item[\textup{(iii)}] El anillo $A/\mathfrak p$ es integralmente cerrado.
\end{enumerate}
En estas condiciones, para todo entero $i$ tal que $0\leq i\leq n$, el anillo
cociente $A/(x_1A+\ldots+x_iA)$ es integralmente cerrado y de dimensión
$\dim(A)\setmin i$. En particular, $\mathfrak J$ es primo, igual a $\mathfrak p$,
$A$ es integralmente cerrado y $(x_i)_{1\leq i\leq n}$ es una sucesión
$A$-regular.
\end{corollary}

Procedamos por inducción sobre $n$, pues la proposición es trivial para $n=0$
y se reduce al lema de Hironaka (5.12.8) para $n=1$. Se puede suponer, por
tanto, que $n\geq2$. Sea $\mathfrak J'=x_1A+\ldots+x_{n-1}A$, y sea
$\mathfrak q$ un ideal primo minimal entre los que contienen $\mathfrak J'$;
se tiene $\operatorname{ht}(\mathfrak q)\leq n\setmin 1$ (\textbf 0, 16.3.1); demostremos que
$\mathfrak q\subset\mathfrak p$. En efecto, si $\mathfrak p'$ es minimal
entre los ideales primos de $A$ que contienen $\mathfrak q+x_nA$,
$\mathfrak p'/\mathfrak q$ es de altura $1$ en $A/\mathfrak q$ por el
Hauptidealsatz (\textbf 0, 16.3.2), por lo que $\mathfrak p'$ es de altura $\leq n$,
pues $A$ es catenario (\textbf 0, 16.1.4); pero, como contiene $\mathfrak J$, es
necesariamente igual a $\mathfrak p$, de donde
$\mathfrak q\subset\mathfrak p$,

\oldpage[IV-2]{131}

y $\mathfrak q$ está inducido sobre $A$ por un ideal primo de
$A_{\mathfrak p}$ minimal entre los que contienen
$\mathfrak J'A_{\mathfrak p}$. Pero, en virtud de la hipótesis (ii) y de (\textbf 0,
17.1.7), $\mathfrak J'A_{\mathfrak p}$ es primo en $A_{\mathfrak p}$, por
tanto
$\mathfrak q=A\cap\mathfrak J'A_{\mathfrak p}$
es el \emph{único} ideal primo asociado no inmerso de $A/\mathfrak J'$ y es de
altura $n\setmin 1$, puesto que $\mathfrak J'A_{\mathfrak p}$ es de altura $n\setmin 1$ y
$\mathfrak q$ de altura $\leq n\setmin 1$. Además,
$\mathfrak qA_{\mathfrak p}=\mathfrak J'A_{\mathfrak p}$ y, como
$A_{\mathfrak q}=(A_{\mathfrak p})_{\mathfrak qA_{\mathfrak p}}$, el ideal
maximal $\mathfrak qA_{\mathfrak q}$ está engendrado por $\mathfrak J'$. Se
ve, pues, que la sucesión $(x_i)_{1\leq i\leq n-1}$ verifica ya las
condiciones (i) y (ii) del enunciado. Demostremos que verifica también (iii).
Consideremos para ello el anillo íntegro $B=A/\mathfrak q$, y sea $t$ la
imagen canónica de $x_n$ en $B$. Es inmediato (cf. (5.6.3, (iv))) que, si $A$
verifica la hipótesis \emph{a)} (resp. \emph{b)}), lo mismo ocurre con $B$; como
$x_n\notin\mathfrak q$, se tiene $t\neq0$; demostremos que $B$ y $t$ verifican
las hipótesis del lema de Hironaka. En efecto, un ideal primo $\mathfrak n$ de
$B$ minimal entre los que contienen $t$ es la imagen de un ideal primo de $A$
minimal entre los que contienen $\mathfrak q+x_nA$, y se ha visto que
$\mathfrak p$ es el único ideal de este tipo; como
$B_{\mathfrak n}=A_{\mathfrak p}/\mathfrak qA_{\mathfrak p}$ y
$\mathfrak nB_{\mathfrak n}
=\mathfrak pA_{\mathfrak p}/\mathfrak qA_{\mathfrak p}$,
el hecho de que $\mathfrak qA_{\mathfrak p}=\mathfrak J'A_{\mathfrak p}$ y
de que $\mathfrak pA_{\mathfrak p}$ esté engendrado por $\mathfrak J$
implica que $\mathfrak nB_{\mathfrak n}$ está engendrado por $t$. Finalmente,
$B/\mathfrak n=A/\mathfrak p$ es integralmente cerrado. La aplicación de
(5.12.8) prueba, pues, que $B=A/\mathfrak q$ es integralmente cerrado y que
$\mathfrak p=\mathfrak q+x_nA$. Utilicemos ahora la hipótesis de inducción,
que muestra que $A/(x_1A+\ldots+x_iA)$ es integralmente cerrado y de dimensión
$\dim(A)\setmin i$ para $0\leq i\leq n\setmin 1$, y que $\mathfrak q=\mathfrak J'$, de
donde $\mathfrak p=\mathfrak J'+x_nA=\mathfrak J$. Se concluye que
$A/\mathfrak J=A/\mathfrak p$ es integralmente cerrado; como
$\dim(A_{\mathfrak p})=n$ y
$\dim(A/\mathfrak p)=\dim(A)\setmin\dim(A_{\mathfrak p})$, puesto que $A$ es
biequidimensional, esto termina la demostración.
